\documentclass[../../infufrgs/ppgc-prop-tese.tex]{subfiles}
\begin{document}


%It is sad to announce that besides a steering method and an in depth analysis of late generalization this work does not entail a new method, or criteria for making models work better, faster or safer. Nonetheless we provide an analysis if interesting, sometimes surprising, behavior of such models and insights. It is a collection of insights about the inner-workings of attention-based models that contribute to a better understanding and appreciation of their emergent phenomena.

This thesis builds on results obtained during the PhD program, including an accepted conference paper and consolidated experimental pipelines. The main completed outcomes are:
%
\begin{enumerate*}[label=(\roman*)]
  \item a set of controlled experiments and analyses on late generalization (grokking-like transitions), consolidated in the paper ``Inducing Grokking with Distribution Shifts'' \cite{carvalho2025inducing_grokking_distribution_shifts};
  \item an implemented prototype of a logic-inspired steering/constraint framework, with preliminary experiments that inform the final methodology and ablation plan; and
  \item an initial set of observations on internal activation/attention patterns under interventions and input perturbations, motivated by sparse-feature diagnostics for steganographic and adversarial perturbations in vision models \cite{carvalho2026interpretable_stego_detection_sae}.
\end{enumerate*}
%
The remaining work consists of consolidating the final steering method, running definitive experiments with baselines and ablations, and integrating all results into a coherent thesis narrative with reproducibility artifacts.

A dedicated prior result that informs the diagnostic track is the SAE-based study on vision models \cite{carvalho2026interpretable_stego_detection_sae}. In that setup, we used pretrained ViT-B/16 as the main backbone and CLIP ViT-B/32 as a complementary model, evaluating clean images against both gradient-based attacks (FGSM/PGD) and low-opacity steganographic text overlays. The experiments showed clear latent-space separation between clean and perturbed samples, and a compact detector using 32 SAE features reached up to 0.90 ROC-AUC for clean-vs-stego discrimination.

For G1, prior evidence supports a boundary-condition interpretation: in controlled synthetic class/subclass settings, distribution-shift induction reliably produced delayed generalization signatures, but transfer to a more realistic clustered-MNIST setup was less consistent \cite{carvalho2025inducing_grokking_distribution_shifts}. This supports using synthetic settings as mechanism-discovery environments while treating real-world transfer as an empirical question rather than an assumption.

For G3, prior evidence is stronger than a single headline metric. The largest activation shifts were concentrated in early-to-mid layers, and competitive detection was obtained with a small attack-sensitive sparse subset rather than full dense activations \cite{carvalho2026interpretable_stego_detection_sae}. This is the main reason the proposal emphasizes layerwise sparse diagnostics plus output-only and dense-feature baselines.

For G2, the current status is feasibility-level maturity: the LoRA+LTN steering architecture and guided decoding logic have been validated conceptually and in pilot form, but the final quantitative matrix remains pending under the frozen protocol and ablation schedule.

In summary, the current evidence establishes feasibility for two pillars of the thesis: delayed-generalization analysis under controlled shift (G1) and sparse-feature diagnostics for perturbation detection (G3). The main remaining uncertainty is whether the full steering-locality track (G2) can consistently improve constraint satisfaction while keeping quality loss and collateral regressions within predefined bounds across seeds and tasks. The definitive experiments that resolve this uncertainty are the frozen G2 protocol runs with complete baseline comparisons and ablations (no LTN, no $T_{\ell\to k}$, and sparsity calibration variants), reported under the unified evaluation matrix in Chapter~\ref{chpt:circuitry}.
\end{document}
